% !TeX root = ../main.tex
\chapter{PENDAHULUAN}

\section{Latar Belakang}

Untuk menelusuri riwayat seekor domba, {\setlength{\fboxsep}{1pt}\colorbox{revisilama}{operator}} {\setlength{\fboxsep}{1pt}\colorbox{revisibaru}{peternak}} perlu mengetahui kapan domba
tersebut ditimbang, pakan apa yang diberikan, serta peristiwa kesehatan dan
reproduksi apa yang pernah dicatat. Setiap catatan perlu memuat identitas domba
yang sama agar {\setlength{\fboxsep}{1pt}\colorbox{revisilama}{operator}} {\setlength{\fboxsep}{1pt}\colorbox{revisibaru}{peternak}} dapat menemukan kembali rangkaian peristiwa tersebut.
Pada ternak sapi, Aisuwarya dkk.~(\citeyear{aisuwarya2019rfid}) mengembangkan
aplikasi untuk mencatat identitas, kesehatan, inseminasi, dan kelahiran ternak.
Pada domba perah, Decandia dkk.~(\citeyear{decandia2026domba}) menggunakan
identifikasi elektronik saat penimbangan untuk mengikuti perubahan bobot
setiap individu dari waktu ke waktu.

Sebagai contoh, {\setlength{\fboxsep}{1pt}\colorbox{revisilama}{operator}} {\setlength{\fboxsep}{1pt}\colorbox{revisibaru}{peternak}} yang membandingkan dua hasil penimbangan seekor
domba mungkin perlu melihat apakah ada perubahan pemberian jenis pakan atau
tindakan kesehatan yang tercatat di antara kedua tanggal penimbangan.
Untuk induk domba, {\setlength{\fboxsep}{1pt}\colorbox{revisilama}{operator}} {\setlength{\fboxsep}{1pt}\colorbox{revisibaru}{peternak}} juga dapat memeriksa apakah ada peristiwa
reproduksi, seperti kelahiran, dalam rentang waktu yang sama. Jika setiap
jenis data disimpan dalam catatan terpisah, {\setlength{\fboxsep}{1pt}\colorbox{revisilama}{operator}} {\setlength{\fboxsep}{1pt}\colorbox{revisibaru}{peternak}} perlu mencocokkan
identitas dan tanggal pada beberapa catatan untuk menyusun urutan peristiwa.
Profil per domba akan memungkinkan {\setlength{\fboxsep}{1pt}\colorbox{revisilama}{operator}} {\setlength{\fboxsep}{1pt}\colorbox{revisibaru}{peternak}} membuka riwayat individu yang
dimaksud dan menelusuri peristiwa yang berdekatan dengan waktu penimbangan.

Ketika penimbangan dan pencatatan dilakukan secara terpisah, {\setlength{\fboxsep}{1pt}\colorbox{revisilama}{operator}} {\setlength{\fboxsep}{1pt}\colorbox{revisibaru}{peternak}} harus
mencocokkan angka pada timbangan dengan identitas ternak sebelum menyalinnya.
Wardhana dkk.~(\citeyear{wardhana2026pemantauan}) melaporkan bahwa peternak
kambing di kandang komunal yang mereka teliti menggunakan kertas atau
mengandalkan ingatan untuk mencatat data ternak. Jika identitas atau waktu
penimbangan terlewat saat pencatatan, {\setlength{\fboxsep}{1pt}\colorbox{revisilama}{operator}} {\setlength{\fboxsep}{1pt}\colorbox{revisibaru}{peternak}} dapat kesulitan menentukan
ternak dan penimbangan yang dimaksud ketika membaca kembali catatan tersebut.
Pada peternakan domba mitra, penelitian ini akan mengamati cara {\setlength{\fboxsep}{1pt}\colorbox{revisilama}{operator}} {\setlength{\fboxsep}{1pt}\colorbox{revisibaru}{peternak}}
mengenali ternak, mencatat peristiwa pada setiap domba, dan menemukan kembali
riwayatnya. Hasil pengamatan akan digunakan sebagai dasar perbandingan dengan
alur pencatatan melalui aplikasi.

Berdasarkan diskusi awal dengan peternak mitra, sebagian informasi mengenai
domba belum dicatat secara rutin. Menurut peternak, ketersediaan riwayat per
domba dapat membantunya menjelaskan kondisi ternak kepada calon pembeli.
Sebagai contoh, ketika calon pembeli menanyakan bobot seekor domba, peternak
dapat merujuk pada hasil penimbangan beserta tanggal pengukurannya.%tambahkan contoh jenis pakan dan riwayat vaksin, karena itu. 

Jika tersedia dalam riwayat pemeliharaan domba yang sama, peternak juga dapat
menunjukkan jenis pakan yang pernah diberikan serta tanggal dan jenis vaksin
yang tercantum dalam catatan vaksinasi.

Bobot juga menjadi pertimbangan dalam pembelian domba untuk potong harian di
Kabupaten Garut. Firdausa dkk.~(\citeyear{firdausa2026preferensi}) melaporkan
bahwa responden pada segmen tersebut paling menyukai kisaran 25--35~kg di
antara pilihan bobot yang diteliti. Kebutuhan yang disampaikan mitra menjadi
salah satu alasan untuk menghimpun catatan pemeliharaan dan hasil penimbangan
pada profil domba yang sama.

Integrasi timbangan dengan aplikasi berbasis \textit{Internet of Things}
(IoT) memungkinkan perangkat mengirim hasil pengukuran langsung ke sistem
pencatatan. Amrulloh dkk.~(\citeyear{amrulloh2024integrasi}) menerapkan
layanan web untuk mengirim hasil timbangan kambing ke platform yang dapat
diakses melalui aplikasi web dan seluler. Wardhana
dkk.~(\citeyear{wardhana2026pemantauan}) menggunakan identifikasi frekuensi
radio (\textit{radio frequency identification}, RFID) untuk mengaitkan hasil
penimbangan dengan identitas kambing. Pada aplikasi yang sama, peternak dapat
memasukkan catatan pemberian pakan dan kesehatan. Dengan rancangan seperti
ini, sebuah profil ternak dapat menghimpun hasil pengukuran dari perangkat
dan catatan yang dimasukkan oleh {\setlength{\fboxsep}{1pt}\colorbox{revisilama}{operator}} {\setlength{\fboxsep}{1pt}\colorbox{revisibaru}{peternak}}.

Riwayat yang terkumpul pada suatu profil tetap diperlukan ketika tag yang
biasa digunakan untuk mengenali ternak hilang. Aisuwarya
dkk.~(\citeyear{aisuwarya2019rfid}) menyediakan fungsi penggantian tag yang
hilang pada aplikasi pencatatan sapi. Pada aplikasi domba yang diusulkan,
{\setlength{\fboxsep}{1pt}\colorbox{revisilama}{operator}} {\setlength{\fboxsep}{1pt}\colorbox{revisibaru}{peternak}} akan membuka profil domba yang telah dikenali, memilih penggantian
tag, lalu mendekatkan tag baru ke pembaca setelah aplikasi menyatakan
timbangan siap. Setelah {\setlength{\fboxsep}{1pt}\colorbox{revisilama}{operator}} {\setlength{\fboxsep}{1pt}\colorbox{revisibaru}{peternak}} menyimpan penggantian tersebut, hasil
penimbangan berikutnya dengan tag baru dapat dicatat pada profil yang juga
memuat riwayat penimbangan dengan tag lama.

Alur penggantian tag dirancang agar {\setlength{\fboxsep}{1pt}\colorbox{revisilama}{operator}} {\setlength{\fboxsep}{1pt}\colorbox{revisibaru}{peternak}} dapat menjalankannya tanpa %add changing rfid without iot knowledge example
memerlukan pengetahuan teknis tentang IoT. Sebagai contoh, jika timbangan
sedang digunakan untuk menimbang, aplikasi akan memberi tahu {\setlength{\fboxsep}{1pt}\colorbox{revisilama}{operator}} {\setlength{\fboxsep}{1pt}\colorbox{revisibaru}{peternak}}
bahwa pemindaian tag baru dapat dicoba setelah penimbangan selesai. Ketika
pemindaian berhasil, aplikasi akan menampilkan kode tag agar {\setlength{\fboxsep}{1pt}\colorbox{revisilama}{operator}} {\setlength{\fboxsep}{1pt}\colorbox{revisibaru}{peternak}} dapat
memeriksanya sebelum menyimpan. Dengan petunjuk ini, {\setlength{\fboxsep}{1pt}\colorbox{revisilama}{operator}} {\setlength{\fboxsep}{1pt}\colorbox{revisibaru}{peternak}} dapat
menentukan tindakan berikutnya dari layar aplikasi, sementara sistem
menangani komunikasi dengan timbangan dan memperbarui penugasan tag.

Penggantian tag juga perlu memperhitungkan hasil yang sudah direkam oleh
timbangan tetapi masih menunggu untuk diunggah. Dalam skenario yang akan
diuji, seekor domba ditimbang menggunakan
tag lama saat koneksi terputus. Setelah koneksi pulih, {\setlength{\fboxsep}{1pt}\colorbox{revisilama}{operator}} {\setlength{\fboxsep}{1pt}\colorbox{revisibaru}{peternak}} mengganti tag
pada profil domba tersebut sebelum perangkat selesai mengunggah hasil
penimbangan. Jika sistem hanya memeriksa penugasan tag yang aktif saat data
diterima, hasil itu dapat dikaitkan dengan domba yang keliru ketika tag lama
digunakan kembali. Pengiriman ulang hasil yang sama juga dapat menghasilkan
catatan ganda apabila sistem memperlakukannya sebagai penimbangan baru.

\begin{sesudahrevisi}
Pengembangan sistem penimbangan dan pencatatan domba ini merupakan proyek
\textit{capstone} yang dikerjakan oleh dua mahasiswa. Rekan satu tim akan
mengembangkan timbangan, termasuk pembacaan tag RFID, penentuan bobot stabil
pada ESP32, dan penyimpanan data di perangkat saat koneksi terputus.
Dalam proyek ini, penulis bertanggung jawab atas aplikasi seluler, basis data,
serta layanan yang menerima dan mencocokkan hasil penimbangan dengan riwayat
penugasan tag. Kedua anggota tim akan menyepakati format data yang dikirim
timbangan dan menguji alurnya hingga hasil yang memenuhi syarat tercatat pada
profil domba yang sesuai.
\end{sesudahrevisi}

Penelitian ini mengusulkan aplikasi seluler yang menyatukan catatan {\setlength{\fboxsep}{1pt}\colorbox{revisilama}{operator}} {\setlength{\fboxsep}{1pt}\colorbox{revisibaru}{peternak}}
dan hasil dari timbangan IoT pada profil domba dengan identitas tetap.
{\setlength{\fboxsep}{1pt}\colorbox{revisilama}{Operator}} {\setlength{\fboxsep}{1pt}\colorbox{revisibaru}{Peternak}} akan menggunakan profil tersebut untuk mencatat dan menelusuri
riwayat kesehatan, pemberian pakan, dan reproduksi. Untuk menjaga
kesinambungan riwayat bobot, sistem akan menyimpan masa berlaku setiap
penugasan tag RFID. Ketika data tiba, layanan penerimaan akan mencocokkan
kode tag dan waktu pengukuran dengan riwayat penugasan serta memeriksa apakah
penimbangan yang sama sudah tersimpan. Proses rekonsiliasi ini memungkinkan
data dari tag lama tetap tersedia untuk penelusuran setelah penggantian tag.

Penggunaan timbangan sebagai sumber data bagi aplikasi berkaitan dengan
bidang lingkungan cerdas dalam draf peta jalan penelitian Departemen Ilmu
Komputer dan Elektronika FMIPA UGM 2023--2027, yang mencakup IoT dan
jaringan sensor \citep[hlm.~22]{dike2023petajalan}. Pada penelitian ini,
pengujian akan mengikuti satu domba dari penimbangan dengan tag lama hingga
penimbangan berikutnya dengan tag baru. Pemeriksaan akan mencakup profil
tempat bobot baru tercatat, keberadaan data dari tag lama setelah unggahan
selesai, dan jumlah catatan bobot setelah perangkat mengirim ulang data yang
sama.

\section{Rumusan Masalah}

Aplikasi pencatatan domba perlu menghubungkan catatan {\setlength{\fboxsep}{1pt}\colorbox{revisilama}{operator}} {\setlength{\fboxsep}{1pt}\colorbox{revisibaru}{peternak}} dan hasil
penimbangan dengan profil domba yang sesuai. Penggantian tag RFID dan
perbedaan antara waktu pengukuran dan waktu unggah perlu diperhitungkan agar
riwayat setiap domba tetap dapat ditelusuri. Berdasarkan kebutuhan tersebut,
penelitian ini merumuskan empat pertanyaan berikut.

\begin{enumerate}
  \setlength{\emergencystretch}{1.5em}
  \item Bagaimana mengembangkan aplikasi seluler yang memungkinkan {\setlength{\fboxsep}{1pt}\colorbox{revisilama}{operator}} {\setlength{\fboxsep}{1pt}\colorbox{revisibaru}{peternak}}
        mengelola data profil domba serta mencatat dan menelusuri riwayat
        kesehatan, pemberian pakan, dan reproduksi?
  \item Bagaimana mengintegrasikan timbangan berbasis IoT dengan aplikasi
        seluler agar hasil penimbangan tercatat pada profil domba yang sesuai
        dan riwayat bobotnya dapat ditelusuri oleh {\setlength{\fboxsep}{1pt}\colorbox{revisilama}{operator}} {\setlength{\fboxsep}{1pt}\colorbox{revisibaru}{peternak}}?
  \item Bagaimana merancang penyimpanan identitas domba dan riwayat penugasan
        tag RFID agar {\setlength{\fboxsep}{1pt}\colorbox{revisilama}{operator}} {\setlength{\fboxsep}{1pt}\colorbox{revisibaru}{peternak}} dapat mengganti tag dan tetap menelusuri
        riwayat bobot melalui profil domba yang sama?
  \item Bagaimana merancang rekonsiliasi data penimbangan agar data yang
        selesai diunggah setelah penggantian tag RFID tetap tersimpan untuk
        penelusuran dan pengiriman ulang data yang sama tidak membuat catatan
        bobot ganda?
\end{enumerate}

\section{Batasan Masalah}%add briding paragraf like in sudais

\begin{sesudahrevisi}
Penelitian ini berfokus pada alur pencatatan yang digunakan {\setlength{\fboxsep}{1pt}\colorbox{revisilama}{operator}} {\setlength{\fboxsep}{1pt}\colorbox{revisibaru}{peternak}}, mulai
dari membuka profil domba hingga menelusuri hasil penimbangannya. Karena
hasil timbang berasal dari perangkat yang dikembangkan bersama dalam proyek
\textit{capstone}, ruang lingkup penelitian perlu memperjelas peran timbangan
dan perangkat lunak pada setiap tahap. Batasan berikut menetapkan objek
penelitian, kondisi pencatatan bobot, serta skenario penggantian tag dan
pengiriman data yang akan diuji.
\end{sesudahrevisi}

\begin{enumerate}
  \item Objek penelitian adalah domba. Alur penggunaan mencakup {\setlength{\fboxsep}{1pt}\colorbox{revisilama}{operator}} {\setlength{\fboxsep}{1pt}\colorbox{revisibaru}{peternak}} yang
        menimbang domba, mengelola profil dan penugasan tag RFID, serta melihat
        riwayat bobot melalui aplikasi seluler. Platform jual beli, alur pembeli,
        dan transaksi merupakan pengembangan lanjutan.
  \item Deteksi bobot stabil dilakukan pada ESP32. Lingkup pengembangan
        perangkat lunak mencakup integrasi timbangan, penerimaan dan rekonsiliasi
        unggahan, basis data, serta aplikasi seluler.
  \item Saat koneksi terputus, timbangan menyimpan data setiap percobaan
        penimbangan, yang disebut satu lintasan, secara lokal untuk diunggah
        setelah koneksi pulih. Pemindaian tag pengganti melalui aplikasi
        memerlukan koneksi antara timbangan dan layanan sisi server.
  \item Hasil penimbangan dapat masuk ke riwayat resmi jika data lengkap, hasil
        ukur valid, dan tag yang terbaca masih aktif serta memiliki tepat satu
        penugasan yang cocok pada waktu pengukuran yang tepercaya. Lintasan
        yang selesai diunggah setelah penugasan tag terkait berakhir tetap
        disimpan untuk penelusuran di luar riwayat bobot resmi. Aturan
        penyimpanan ini juga berlaku untuk lintasan dengan identitas meragukan
        atau waktu pengukuran yang tidak tepercaya. Konfirmasi identitas melalui
        aplikasi merupakan pengembangan lanjutan.
  \item Pengujian bobot menggunakan timbangan acuan. Skenario pengujian
        mencakup pembacaan yang tidak stabil dan pembacaan yang tampak stabil
        ketika domba belum sepenuhnya berada di atas platform. Pengujian pencatatan
        meliputi penggantian tag, konflik penugasan tag, waktu pengukuran yang
        meragukan, unggahan sebagian, dan pengiriman ulang.
\end{enumerate}

\section{Tujuan Penelitian}

Penelitian ini bertujuan mengembangkan aplikasi seluler yang memungkinkan
{\setlength{\fboxsep}{1pt}\colorbox{revisilama}{operator}} {\setlength{\fboxsep}{1pt}\colorbox{revisibaru}{peternak}} mengelola profil domba serta mencatat dan menelusuri riwayat
kesehatan, pemberian pakan, dan reproduksi. Aplikasi akan diintegrasikan
dengan timbangan berbasis IoT agar hasil penimbangan dapat dicatat pada
profil domba yang sesuai. Penelitian ini juga akan merancang penyimpanan
identitas tetap dan riwayat penugasan tag RFID, serta mengembangkan layanan
untuk menangani unggahan tertunda dan pengiriman ulang. Pengujian akan
memeriksa fungsi pencatatan dan penelusuran riwayat, kesinambungan profil
setelah penggantian tag, keterlacakan data yang terlambat diunggah, serta
pencegahan catatan bobot ganda saat data yang sama dikirim ulang.

\section{Manfaat Penelitian}

Secara praktis, aplikasi ini diharapkan membantu {\setlength{\fboxsep}{1pt}\colorbox{revisilama}{operator}} {\setlength{\fboxsep}{1pt}\colorbox{revisibaru}{peternak}} menemukan riwayat
bobot, kesehatan, pemberian pakan, dan reproduksi melalui satu profil domba
serta menelusuri kembali riwayat bobot setelah tag RFID diganti. Secara
akademik, rancangan dan hasil pengujian penelitian ini diharapkan dapat
menjadi rujukan bagi pengembangan sistem pencatatan ternak berbasis IoT yang
menggunakan identitas tetap dan riwayat penugasan tag, serta bagi pengujian
sistem yang menerima unggahan penimbangan tertunda atau berulang.
